\section{Results}
\label{sec:results}

\subsection{The prefill arc (v1--v5, T4)}
\label{sec:results-prefill}
The first five kernels rebuild FlashAttention for prefill and establish the
methodology; each pure-roofline miss teaches something the traffic model cannot
see. \textbf{v1} (naive, three-pass FP32) is HBM-bound in principle
($\ai \approx 0.25$) but runs \emph{faster} than its cache-free floor at mid-$N$
(85.6\,ms vs a 109\,ms floor at $N{=}2048$, $d{=}64$) because the T4's 4\,MB L2
absorbs the redundant operand reads; only at $N{=}8192$, when L2 overflows, does
it converge to the floor. \textbf{v2} (shared-memory tiling) is
$1.3$--$3.2\times$ faster than v1, but \texttt{ncu} shows the win is
\emph{compute/scheduling, not bandwidth}: DRAM traffic is essentially unchanged
($1.02\times$ at $N{=}8192$), because the L2 had already served the operands and
because $\sim$99\% of DRAM traffic is the materialized score matrix $S$, which is
byte-identical in v1 and v2. The roofline predicted a $\sim$30$\times$ traffic
cut; the measured cut was zero. \textbf{v3} (online softmax) eliminates $S$ from
HBM entirely, dropping peak memory $125\times$ (2164\,MB $\to$ 17\,MB), yet is
$3$--$7\times$ \emph{slower} than v2, proving $S$ was never the bottleneck on a
machine that was never bandwidth-bound; the kernel sits $151\times$ above its MMA
floor, occupancy/latency-bound. \textbf{v4} (fused single-pass, warp-per-row)
recovers the schedule ($1.7$--$2.6\times$ over v2, $7.5$--$15\times$ over v3) and
closes the distance to the floor from $151\times$ to $\sim$18$\times$; the
residual gap is FMA under-utilization, since warp-per-row scoring is GEMV-shaped,
a per-key shuffle reduction dominating the two FMAs. \textbf{v5} (tensor-core WMMA)
is predicted to drop the floor $8\times$ but its prefill benchmark was
\emph{never executed}; I report it as a prediction only and never as a measured
result (Section~\ref{sec:threats}).

\subsection{The decode arc (v6--v12)}
\label{sec:results-decode}
The decode arc is the contribution. All decode readings use the counter-free
proxy of Section~\ref{sec:measure}.

\paragraph{Split-KV and paging (v6--v7)}
Decode with $N_q{=}1$ leaves a single-query kernel with one CTA and idle SMs.
Split-KV partitions the KV across blocks, each emitting an unnormalized partial
$(O, m, \ell)$ that an LSE merge recombines. \textbf{v6} beats the naive
$N_q{=}1$ loop by $5.7$--$8.2\times$ and torch SDPA by $1.5$--$3.3\times$, but
reaches only \textbf{9--15\% of HBM} ($6.5$--$11.6\times$ above the bandwidth
floor): the limiter is occupancy/launch, not bandwidth. \textbf{v7} adds a paged
block-table gather (the vLLM cache layout) and, critically, a \emph{batch
sweep} that \emph{refutes} the predicted occupancy$\to$bandwidth crossover: \%HBM
stays flat at $9.4$--$12.4\%$ from $BH{=}8$ to $BH{=}512$
(Table~\ref{tab:v7batch}). Code inspection confirms the cause: shared memory
caps residency at 2 blocks/SM, and at $N_q{=}1$ only one of eight warps computes,
so batch merely adds waves. Decode is \emph{per-CTA-bound at every batch size},
not grid-occupancy-bound, which reorders the roadmap: per-CTA efficiency
(GEMV$\to$GEMM) before bytes.

\begin{table}[t]
  \centering
  \caption{v7 batch sweep at $N_k{=}8192$ (T4). \%HBM is flat across a
  $64\times$ batch range: no occupancy$\to$bandwidth crossover. v7 also loses to
  SDPA at $B{\ge}8$ (SDPA amortizes launch; v7 is already SM-saturated at $B{=}1$).}
  \label{tab:v7batch}
  \small
  \begin{tabular}{rrrrr}
    \toprule
    $B$ & $BH$ & \%HBM $d{=}64$ & \%HBM $d{=}128$ & vs SDPA $d{=}64$ \\
    \midrule
    1  & 8   & 9.8\%  & 12.2\% & 2.65$\times$ \\
    8  & 64  & 9.4\%  & 11.4\% & 0.56$\times$ \\
    16 & 128 & 9.4\%  & 11.4\% & 0.56$\times$ \\
    32 & 256 & 9.5\%  & 11.6\% & 0.51$\times$ \\
    64 & 512 & 10.3\% & 12.4\% & 0.50$\times$ \\
    \bottomrule
  \end{tabular}
\end{table}

\paragraph{GQA $M$-packing and score-stationary (v8, v8.7)}
Packing the $G$ query heads of a GQA group into the CTA's $M$ dimension activates
$G$ warps (not one) and reads the shared KV once, raising \ai\ from $2/b$ to
$2G/b$. On plain CUDA cores this \textbf{v8 Cut~1} tracks $\sim G\times$ up to
$G{=}8$ (measured $8.59\times$ at $d{=}64$, $8.71\times$ at $d{=}128$ versus
no-packing) and beats SDPA $6$--$10\times$ at \emph{every} batch $B{=}1$--$64$,
reclaiming the serving regime v7 had lost. \%HBM stays $\le 11\%$: the roofline
got the speedup \emph{scaling} right (a first partial win in six steps) while the
kernel never approaches the absolute floor. Tensor cores are the \emph{wrong}
tool here: a Turing WMMA variant (Cut~2) is $2.1$--$3.3\times$ \emph{slower} than
Cut~1 even at $G{=}16$ (a full WMMA tile), pinned by the opaque-fragment
staging tax, not compute. Two further single-variable ablations, double-buffered
KV (v8.5) and key-ILP/occupancy (v8.6), were both null, isolating the floor as
the per-row serial online-softmax recurrence. \textbf{v8.7} removes it: a
score-stationary relayout (lane $=$ key) eliminates the per-key cross-lane
reduction and shortens the recurrence $32\times$. It is the first inner-loop
lever to move the clock-robust \%HBM (up $2$--$3$ points, to $\sim$10--12\%) and
beats Cut~1 by $1.1$--$1.6\times$ at matched clock. This confirms the recurrence
was a real, \emph{removable} floor component, while a residual per-CTA ceiling
remains.

\paragraph{Cross-precision KV (v9--v10)}
\textbf{v9} stores the KV cache in FP8 E4M3 (1 byte), dequantized fused per tile.
The pre-registered ``capacity-only'' prediction is \emph{refuted}: FP8 delivers a
real, clock-matched $\sim$1.3$\times$ decode-latency win (median across $G$;
$0.96$--$1.52\times$), and the win \emph{shrinks as $G$ grows}
($1.37\times$ at $G{=}2 \to 0.96\times$ at $G{=}32$, $d{=}64$). That shrinkage is
the fingerprint of an L2$\to$SM \emph{load} effect that $M$-packing amortizes
away, not an HBM effect. \%HBM \emph{drops} versus FP16 and never flips the limiter; per-tensor
E4M3 RMSE versus FP16 KV is $\sim$6--7$\times10^{-4}$. \textbf{v10} stores NVFP4
(packed E2M1 nibble $+$ per-16 micro-scale, 0.5625\,B/element): measured capacity
is $3.56\times$ versus FP16 and $1.78\times$ versus FP8. NVFP4 RMSE is
$\sim$2.4$\times10^{-3}$, about $4\times$ FP8's (E2M1's single mantissa bit, not
a bug), and there is \emph{no} latency win (\%HBM falls to $\sim$3\%). An
asymmetric per-channel-K / per-token-V recipe recovers accuracy only where real
(GPT-2) KV has channel-outlier structure and is layer-conditional, so I ship the
robust block-16 default.

\paragraph{v9 Task~1: the confound-free per-CTA verdict}
The ``$\sim$10\% HBM, per-CTA-bound'' reading recurs from v6, but until v9~Task~1
it was confounded: the micro-bench KV fit in L2, clocks were unlocked, and $N_k$
never passed L2. On a \emph{root} T4 I removed all three at once (clocks locked
at 1590\,MHz, L2 flushed, $N_k \to$ 128K so the working set is 537\,MB $\gg$ 4\,MB
L2) with \texttt{ncu} finally live, the first counter-validated measurement in
the project. Achieved \%HBM plateaus at $\sim$28--29\% under full occupancy (never
near the $\sim$70\% achievable ceiling), with \emph{no bandwidth knee} at the L2
crossing; \texttt{ncu} reports L2 hit-rate 1.1\% (data genuinely from HBM) yet
DRAM throughput only 12.85\%. The verdict is earned: decode is
\emph{per-CTA / low-MLP latency-bound}, not bandwidth-bound, and the counter-free
proxy is validated (13.8\% vs 12.85\%). \textbf{Bytes are not the decode wall}:
FP8/NVFP4 are a capacity$+$accuracy play, not a decode-latency play.

\paragraph{MLA decode: the shape change (v11--v12)}
With bytes exhausted and occupancy capped, the only lever left is the
\emph{shape}. \textbf{v11} implements MLA on CUDA cores: all $h_q{=}128$ heads
share one latent, so $M{=}128$ by construction and \ai\ jumps to $\approx$235.
Pure roofline calls this the arc's first compute-flip; the per-CTA-corrected
layer predicts it \emph{stays} per-CTA because a warp-per-head kernel packs
$M{=}128$ as $16$ blocks $\times 8$ warps, not one GEMM. Measurement confirms the
corrected layer on \emph{both} T4 and B300: per-CTA-bound to $N_k{=}2$M past a
$5\times$ L2 overflow (\%HBM $\approx 0$), at $0.75\tflops$ ($<$1\% of either
compute peak). The headline is a $\sim$4$\times$ gap: torch dense-MQA, which
routes $M{=}128$ into cuBLAS tensor-core GEMMs, beats my CUDA-core kernel
$\sim$4$\times$ on B300. This is a clean result either way: the MLA \emph{shape}
is correct and tensor-core-friendly, but the CUDA-core default is the wrong tool.
(I attribute the gap to cuBLAS exploiting the tensor cores; the torch baseline
itself was not profiled, so the ``$M{=}128$-as-one-TC-GEMM'' benefit is shown
structurally and externally~\cite{gla}, not by a torch-side measurement.)
Capacity is $202\times$ versus MHA / $12.6\times$ versus GQA-8, and the latent
absorption identity (absorbed kernel $\equiv$ explicit per-head materialization)
is verified in the correctness suite.

\textbf{v12} changes only the engine. On the B300 box I found that CUTLASS 4.6
example~77 already ships the production \texttt{tcgen05} MLA decode kernel
(\texttt{Sm100FmhaMlaKernelTmaWarpspecialized}: $M{=}128$ by
\texttt{static\_assert}, FP8-native, paged $+$ split-KV $+$ LSE), so v12's
deliverable became characterizing that canonical kernel directly rather than
reinventing NVIDIA's scheduler. Figure~\ref{fig:regime} and Table~\ref{tab:v12}
give the result: at realistic single-request decode ($B{=}1$, $K{\le}32$K) even
the production kernel runs at \textbf{1--2\% of FP8 peak}, exactly the
pre-registered per-CTA/pipeline-depth floor; but throughput \emph{scales with
total work $B\times K$} to \textbf{1785\tflops\ $=$ 35.7\% of FP8 peak / 46.3\% of
HBM} at $B{=}64$, $K{=}524{,}288$. FP16 tracks FP8 (not compute-bound), and
\texttt{--verbose} confirms 174\,KB shared memory $\to \sim$1 CTA/SM (the measured
capacity pressure). The transition is smooth in $B\times K$ with no hard knee.

\begin{figure}[t]
  \centering
  \includegraphics[width=\columnwidth]{v12-throughput-regime.pdf}
  \caption{v12 (production \texttt{tcgen05} MLA decode, B300): FP8 throughput
  versus total work $B\times K$. Single-stream decode ($B{=}1$, moderate $K$)
  sits at 1--2\% of the 5\,PF FP8 peak; throughput climbs to 35.7\% peak / 46.3\%
  HBM as $B\times K$ grows. The floor is work-starvation, not a fixed wall.}
  \label{fig:regime}
\end{figure}

\begin{table}[t]
  \centering
  \caption{v12 FP8 decode throughput vs work (B300, CUTLASS ex77, \texttt{sm\_103a}).
  Peaks: FP8 5\,PF, HBM 8\tbs.}
  \label{tab:v12}
  \small
  \begin{tabular}{rrrrr}
    \toprule
    $B$ & $K$ & \tflops & \% FP8 peak & \% HBM \\
    \midrule
    1  & 256    & 2.4    & 0.05\% & 0.15\% \\
    1  & 8192   & 51.4   & 1.0\%  & 1.4\%  \\
    1  & 32768  & 112.5  & 2.2\%  & 2.9\%  \\
    1  & 524288 & 790.6  & 15.8\% & 20.5\% \\
    64 & 8192   & 944.8  & 18.9\% & 25.5\% \\
    64 & 131072 & 1717.3 & 34.3\% & 44.5\% \\
    64 & 524288 & \textbf{1785.3} & \textbf{35.7\%} & \textbf{46.3\%} \\
    \bottomrule
  \end{tabular}
\end{table}

\begin{figure}[t]
  \centering
  \includegraphics[width=\columnwidth]{v12-work-starvation-correction.pdf}
  \caption{The correction. v6--v11 CUDA-core decode is flat near the per-CTA
  floor (v11 MLA: $0.75\tflops$ regardless of $B\times K$). The \texttt{tcgen05}
  engine (v12) spans $2.4 \to 1785\tflops$ as work grows, so ``per-CTA-bound
  \emph{forever}'' was a small-$B\times K$ artifact; the wall is work-starvation
  and the engine is the lever.}
  \label{fig:starvation}
\end{figure}

This \emph{corrects} the hedged v6--v11 reading. ``Per-CTA-bound \emph{forever}''
was a small-$B\times K$ artifact of the CUDA-core kernels: v11's warp-per-head
GEMV is flat at $0.75\tflops$, whereas the same-shape \texttt{tcgen05} kernel
spans $2.4 \to 1785\tflops$ ($\sim$70$\times$ v11 at $B{=}1$/$K{=}8$K,
$\sim$1000$\times$ at $K{=}524$K). The per-CTA wall is work-starvation; the right
\emph{engine} converts work into throughput (Figure~\ref{fig:starvation}). I
measured NVIDIA's canonical kernel; I complement, not beat it. (Honest
deviation: the pre-registered v12 row assumed NVFP4 storage, \ai\ 835; ex77 is
FP8-storage, \ai\ 470 $<$ the FP8 ridge 625, so pure roofline predicts HBM-bound,
and at scale the kernel indeed leans to the HBM roof, 46.3\% HBM above
its 35.7\% compute fraction. NVFP4-storage and the native-FP4 \emph{compute} arm
remain unmeasured; see Section~\ref{sec:stagea}.)

\subsection{The cross-architecture per-CTA ceiling}
\label{sec:results-crossarch}
The strongest evidence that the decode floor is a \emph{shape} property, not a
hardware limit, is architecture-independence. At $B{=}1$ decode the CUDA-core
kernels achieve $\sim$40\gbs\ on T4, B200, \emph{and} B300
(Table~\ref{tab:crossarch}): 11--14\% of the T4's 320\gbs\ but only $\sim$0.5\%
of Blackwell's 8\tbs. If decode were bandwidth-bound the \emph{percentage} would
be constant; instead the \emph{absolute} $\sim$40\gbs\ ceiling holds across a
$25\times$ span of peak bandwidth, which is the per-CTA/low-MLP latency signature.
On B300 the CUDA-core MLA kernel confirms this to $N_k{=}2$M (a $5\times$ L2
overflow), with \%HBM flat near zero throughout.

\begin{table}[t]
  \centering
  \caption{The architecture-independent $\sim$40\gbs\ per-CTA ceiling at $B{=}1$
  decode. The absolute achieved bandwidth, not the percentage, is invariant
  across a $25\times$ span of peak HBM.}
  \label{tab:crossarch}
  \small
  \begin{tabular}{lrrl}
    \toprule
    Arch & HBM peak & achieved @ $B{=}1$ & \% of peak \\
    \midrule
    T4 (\texttt{sm\_75})   & 320\gbs & $\sim$36--46\gbs & 11--14\% \\
    B200 (\texttt{sm\_100}) & 8\tbs   & $\sim$40\gbs      & $\sim$0.5\% \\
    B300 (\texttt{sm\_103}) & 8\tbs   & $\sim$40\gbs      & $\sim$0.5\% \\
    \bottomrule
  \end{tabular}
\end{table}
